\honest{Beyond the Reach of God}{Eliezer Yudkowsky}{2008}

This post is a gloomy thought experiment I made up to smash my own optimism. In ``The Magnitude of His Own Folly'' I traced my folly to a feeling that everything would be all right in the end. Readers who want to keep their biases because the biases keep them happy ``should consider not reading,'' unless they have something to protect. Some count that inner optimism a virtue; others call it necessary for mental health. But we live in the world beyond the reach of God.

I grew up in an Orthodox Jewish family, and I stopped believing in God long ago. I remember my last prayer, on behalf of the boy next door, and what it was like to have a higher authority to take care of what I could not. Take God as seriously as a child does. Even then God need not fetch you a lemonade from the refrigerator. If you said you had built a benevolent superintelligence and I asked for a banana and none appeared, that would not yet disprove you. Parents do not grant every wish. I accept fun-theoretic arguments against always giving people what they want at once. I do not want to become a simple wanting-thing that never has to plan or act. So saying that God does not grant all prayers is not necessarily an evasion; even a Friendly AI might not answer every request.

Still, ``clearly, there exists some threshold of horror awful enough that God will intervene.'' A God who never intervenes, however bad things get, is ``an obvious attempt to avoid falsification.'' Young children really expect to see the dragon in their garage. Even a child can imagine arguing over where the threshold lies, but God will draw the line somewhere. Few loving parents would let their toddler be run over by a car. \nb{Here the post takes a side in the problem of evil. The theist reply, as the Stanford Encyclopedia of Philosophy reports it, is that the goods that would justify God in permitting an evil may lie beyond our ken. On that reply, God's not intervening in the horrors we know of is not a belief built to avoid falsification. William Rowe answered that the reply, if it worked, would undercut inductive reasoning in general.}

The plainest horror God could not tolerate is true death, annihilation of the mind. In the world where God exists, no soul need fear it. Now play Conway's Game of Life, which helps in understanding what a physical law is. It is Turing-complete, so it could hold a sentient being. If you built Conway's Game of Life on your computer and set it to torture a sentient being, God would step in and change your transistors. \nb{Life can compute whatever a universal computer can, as the post says.}

But ask instead: given these initial conditions and these rules, what would the mathematical result be? Not even God can change that answer, unless God can do the logically impossible. In that world each step follows only from the one before. It is not fair. Life may or may not evolve, with no God to guide it. In that what-if world, conscious cows may be eaten by conscious wolves, diseases teach no lessons, and ``the equivalent of Genghis Khan can murder a million people, and laugh,'' grow rich, never be punished and live more happily than most.

If he tortures people to death for days, they may call out to an imagined God. Since the laws contain no prohibition against torture, they are saved only if the right cells happen to be 0 or 1. Anyone who defies the Khan is struck with a sword and dies. The victims could be wholly innocent. Life's rules are very simple: a cell with three living neighbors is alive at the next step, a cell with two stays as it is, and all other cells die. They say nothing about innocence. ``Is this world starting to sound familiar?'' \nb{Paul Draper made this comparison formally in 1989. Draper argued that the facts of pain and pleasure are more probable if neither the nature nor the condition of sentient beings on earth is the work of benevolent or malevolent nonhuman persons than if God exists.}

Belief in a fair universe shows in subtler ways. Would the twentieth century have gone differently if Hitler's parents had made love an hour earlier and a different sperm had fertilized the egg? So many lives turning on one small event seems disproportionate. The Divine Plan ought to make more sense. So, the believer concludes, if Hitler had become an architect, someone else would have taken his role. But nothing in the laws of physics says that big effects need big causes. You can believe in a Divine Plan without God: ``Karl Marx surely did.'' \nb{Karl L\"owith argued in 1949 that philosophies of history up to Marx's were secularised versions of Christian eschatology.} The point of the thought experiment is to lay the God-universe and the Nature-universe side by side, to recognize which thoughts belong to the first. Many atheists still think certain things are not allowed, and would argue that World War II would have happened without Hitler. I call this ``an unreasonable belief'', since ``from my reading, it seems that events were mostly driven by Hitler's personality''. \nb{That reading agrees with the long-standing general view that Hitler wanted, planned and launched the war.} The main reason to doubt it ``would be refusal to accept'' that the universe could make so little sense. I name no one who doubts it. \nb{The best-known dissent, A.~J.~P. Taylor's in 1961, argued from diplomatic documents that Hitler's goals were the same as those of other German politicians such as Gustav Stresemann. Its critics rejected it on the evidence.}

In the God-universe, God prohibits such things. In ours, whatever physics says will happen, will happen, even in the most extreme cases. Reading William Shirer's \textsc{The Rise and Fall of the Third Reich}, with its account of how Shirer and others disbelieved the full scope of Nazi atrocities, I was ``horrified, but not at all disbelieving,'' because I already knew there was no protection against it. I once believed that the extinction of humanity was not allowed, and other rationalists ``may yet have things they trust'': positive-sum games, democracy, technology.

These are held sacred: they cannot lead to anything really bad, or not permanently without a silver lining. History cannot turn from a positive-sum trend to a negative-sum one; modern liberal democracies will never legalize torture; no Black Swan technology will do more harm than all the good so far. Their ``clever arguments'' come from ``a much deeper belief that such things are not allowed.'' \nb{The only case the post gives is the author's own former belief.} Until you can picture a lawful universe where such things happen, you ``aren't yet ready to argue probabilities.''

Could sentient beings have died absolutely for millions of years, as part of no grand plan and teaching no lesson, so that a trick as simple as vitrifying people in liquid nitrogen could save them, and a ten-second rejection of the idea could destroy someone? Could a programmer who signs some papers and buys life insurance continue into the far future while Einstein rots in a grave? God would not allow anything so disproportionate. Secular rationalizations can do the same work, so it helps to imagine a benevolent God who enforces a minimum of fairness, and then to imagine that you yourself live outside that God's reach, in a world of pure mathematics where anything can happen. \nb{The book's own discussion of cryonics is ``When Science Can't Help''.} Readers who want happiness above all should not dwell on this, beyond signing up for cryonics, wearing a seat belt, getting health insurance and the other dreary necessities. They might also ``write a check to an existential-risk-mitigation agency now and then.''

This post is for those with something to protect. A twelfth-century peasant could do nothing against annihilation. Nature offers ``Absolute, utter, exceptionless neutrality.'' Challenges are not always fair. A lethal penalty means death, for people and for planets. Knowing this will not always save you. If you think a rationalist who understands the mess must surely find a way out, then you trust rationality. Some commenter will list the reasons it is lovely to live in a neutral universe: life may be a little dark, but not past a certain point without a silver lining.

Still, a few hopeful words. We cannot change fundamental physics, but at a higher level we might build guardrails and padding within our future light cone, the only world it does any good to care about; everything outside it we can treat as a ``generalized past.'' Someday children might reliably be sheltered: they might burn a finger but never be run over by cars. A superintelligence, which could think a trillion thoughts without a misstep, would find the raw universe only another problem. But ``building an adult is itself an adult challenge.'' A fairer universe must be reached from this one, where challenges are not calibrated to your skills. Not everyone needs these unpleasant thoughts; buckling a seat belt is not so hard. But anyone who plans to confront an uncalibrated challenge of instant death must not avoid them. I end by asking what rules a child should follow ``to solve an adult problem.''
